\section{Conclusion}\label{sec:conclusion}

For linear elastostatics, the assembled discrete potential energy is an exact label-free
training objective: minimising it is supervised regression in the stiffness norm, which
measures the strain energy of the stress error, and it needs no solve. The same identity
makes the energy an evaluation metric that follows the stress rather than the displacement
and ranks the predictions of one load case without reference solutions, and a failure
detector: the sign of the energy of a prediction tells whether it is worse than the zero
field, and a closed-form rescaling ensures that it is not. In a pre-registered comparison on
three-dimensional tetrahedral meshes, a transformer trained on the energy alone was more
accurate on every metric, in each of three seeds, than the same transformer trained on the
solutions of the same instances with the relative displacement error. In the same run, by a
reading not fixed in advance, a supervised graph network with lower displacement errors had
several times larger von Mises errors, so that the two errors ranked the networks in
opposite orders. These findings replicated on held-out instances. Trained on the
stiffness-norm error of the same solutions, the supervised transformer lowered its median
relative energy gap by \numKdGapRel{} and its median von Mises error by \numKdVmRel{} (on
the seeds' medians, a statistic chosen after the per-instance values of the networks
compared were known), closing its gap to the label-free transformer, whose median relative
energy gap it undercut by a margin that three seeds leave uncertain: what the label-free
objective gains over displacement supervision is the norm, which it provides without a
solve. In two dimensions the label-free network had a larger displacement error than the
supervised one and smaller energy and stress errors; trained on the stiffness-norm error
instead of the displacement error, with the same labels, the supervised network lowered its
von Mises error by more than half, supporting a pre-registered hypothesis that replicated on
held-out instances, and its errors lay lower in the spectrum of the stiffness matrix, as
measured afterwards; the reduction of its relative energy gap did not replicate there. In
three dimensions, a classical Ritz projection of a few hundred unknowns, added after the
original analyses and given the stiffness matrix, was at its highest degree more accurate in
displacement than each network on almost every instance, and less accurate in peak stress
than the two transformers trained in the stiffness norm. A warm start from one label-free
prediction of the earlier code did not shorten unpreconditioned conjugate gradients to a
relative residual of \numKfiveTol{} in two dimensions, and the predictions timed in
\cref{tab:cost} did not shorten it to \numCostCgTol{} in two dimensions or three. In three
dimensions no network transferred to a substantially finer mesh without retraining, although
a label-free rescaling removed most of the label-free network's error there. These limits,
and the criteria that were not met, are reported with the results.
